\documentclass[10pt,a4paper]{amsart}
\usepackage[margin=19mm]{geometry}
\usepackage{amsmath,amssymb,mathtools}
\usepackage[scr=rsfs,cal=euler]{mathalfa}
\usepackage{xfrac}
\usepackage[unicode,hidelinks,bookmarks=false]{hyperref}
\newcommand{\ie}{i.e.\ }
\newcommand{\cf}{cf.\ }
\newcommand{\resp}{respectively\ }
\newcommand{\Subsection}{\subsection}
\providecommand{\nobreakdash}{\nobreak\hbox{-}}
\setlength{\emergencystretch}{2em}
\multlinegap0pt
\usepackage{bm}
\usepackage{bbm}
\usepackage{dsfont}
\usepackage{xspace}
\usepackage{cancel}
\usepackage{mathtools}
\usepackage{amsthm}
\makeatletter
\@ifundefined{proposition}{\newtheorem{proposition}{Proposition}}{}
\makeatother
\newcommand{\hMat}{\mathbf H}
\newcommand{\pVec}{\mathbf p}
\title[Hill-Type Stability Analysis of Periodic Solutions of Fracti]{Hill-Type Stability Analysis of Periodic Solutions of Fractional-Order Differential Equations}
\author{Paul-Erik Haacker, Remco I. Leine, Renu Chaudhary, Kai Diethelm, Andr\'e Schmidt, Safoura Hashemishahraki}
\date{Source: arXiv:2509.24639v3 (CC BY 4.0)}
\begin{document}
\maketitle

\noindent\emph{Source and licence.} Hill-Type Stability Analysis of Periodic Solutions of Fractional-Order Differential Equations. Version arXiv:2509.24639v3 (CC BY 4.0), distributed under the Creative Commons Attribution 4.0 International licence, as recorded in the arXiv licence metadata for this exact version and shown on the abstract page https://arxiv.org/abs/2509.24639v3. Full rights notice: licenses/arxiv-2509.24639-CC-BY-4.0.txt.\par\smallskip

\noindent\textit{Editorial scope.} Counted unit: the Proposition whose source label is \texttt{prop\_ComplexConj} in the article (source0.tex lines 1251--1253, source label \texttt{prop\_ComplexConj}), delivered together with the article's own complete original proof of it (source0.tex lines 1254--1272, one \texttt{proof} environment of 19 lines). No mathematics is written, repaired or re-derived: statement and proof are the article's own text line for line. Required context retained uncounted: eq\_FracHillFourierForm (lines 1087--1089); eq\_solutionFloquetForm (lines 1095--1097); eq\_fractionalHillProblem (lines 1230--1232). Every definition, display and label of those lines is retained unchanged. Omitted from this package and not redistributed: 1--1086, 1090--1094, 1098--1229, 1233--1250, 1273--1722. Nothing inside a retained excerpt is omitted or altered: every retained body is delivered exactly as the article's lines are written, so no omission receipt of any kind accompanies this package. Citations: the retained excerpts carry no citation command at all, so no bibliography entry of the article is transcribed and no reference is redistributed. Reference adaptation: none was needed and none was made; every reference inside the delivered unit resolves inside it, so no unresolved reference or citation is produced. Printed slips kept literal: no printed word, symbol, hypothesis or number of the counted statement or of its proof was altered. Layout and class: the article's own class is \texttt{svproc} with packages including geometry, lipsum, url, xcolor, mathtools, ulem, enumitem, subcaption, array, hyperref, graphicx, epstopdf, tikz, pgfplots; the wrapper is the lane's pinned \texttt{amsart} editorial wrapper at 10pt on A4, which restates the theorem-like environments the retained text needs and adds the article's own macro definitions that the retained text needs (\texttt{\textbackslash hMat}, \texttt{\textbackslash pVec}), with extra wrapper packages (bm, bbm, dsfont, xspace, cancel, mathtools). The delivered numbering is the unit's own. Assets: no retained span contains a figure environment, a graphics-inclusion command, a PDF-page inclusion, a table or a TikZ picture; no image file is copied into this package, the package declares no asset dependency and no image file is redistributed with it. Licence: version arXiv:2509.24639v3 of this article is distributed under the Creative Commons Attribution 4.0 International licence, as recorded in the arXiv licence metadata for this exact version and shown on the abstract page https://arxiv.org/abs/2509.24639v3. A full rights notice is delivered at licenses/arxiv-2509.24639-CC-BY-4.0.txt. Characters: every retained excerpt is pure ASCII, so the delivered engine, XeLaTeX with its default Latin Modern text font, reports no missing character.
				\begin{align}\label{eq_FracHillFourierForm}
					J(t) = \sum_{k=-\infty}^{+\infty}J_k\exp(i\omega k t),\quad p(t) = \sum_{k=-\infty}^{+\infty}p_k\exp(i\omega k t)
				\end{align}

				\begin{align}\label{eq_solutionFloquetForm}
					y(t) = \exp(\lambda t)p(t)
				\end{align}

		\begin{align}\label{eq_fractionalHillProblem}
			0 = \det \hMat_{N}^{(\alpha)}(\lambda).
		\end{align} This truncated fractional Hill problem then allows for computation of $\pVec_N$ as an element of the nullspace of $\hMat_{N}^{(\alpha)}(\lambda)$ and thus for construction of the solution \eqref{eq_solutionFloquetForm}. The equivialent formulation

		\begin{proposition}\label{prop_ComplexConj}
			The solutions of the truncated fractional Hill problem \eqref{eq_fractionalHillProblem} come in complex conjugate pairs $(\lambda,\bar\lambda)$. 
		\end{proposition}

		\begin{proof}
			First note two things:
			\begin{enumerate}
				\item Since $J(t)$ is real, its Fourier coefficients \eqref{eq_FracHillFourierForm} satisfy $\overline{J_k} = J_{-k}$ for all $k$, where $\overline{(\cdot)}$ denotes the element-wise complex conjugate.
				\item The principal root function $\lambda^\alpha = |\lambda|^\alpha \exp(i\alpha \arg \lambda)$ commutes with complex conjugation, i.e. $\overline{(\lambda^\alpha)} = (\overline \lambda)^\alpha$. 
			\end{enumerate}
			For a quadratic matrix $M$, it holds that $\operatorname{rank}M = \operatorname{rank}\overline{M}$ and we may take $M = \hMat_{N}^\alpha(\overline\lambda)$
			\begin{align}
				\operatorname{rank} \hMat_{N}^\alpha(\overline\lambda) = \operatorname{rank}\overline{\hMat_{N}^\alpha(\overline\lambda)}
			\end{align}
			where the latter matrix -- using the two statements above -- reads
			\begin{align}\label{eq_proofComplexConjMat1}
				\overline{\hMat_{N}^\alpha(\overline\lambda)} = \begin{pmatrix}
					J_0 - (\lambda +iN\omega)^\alpha I & \cdots & J_{2N}\\
					\vdots & \ddots & \vdots \\
					J_{-2N} & \cdots & J_0 - (\lambda -iN\omega)^\alpha I
				\end{pmatrix}.
			\end{align} Swapping columns and rows blockwise, we see that matrix \eqref{eq_proofComplexConjMat1} has the same rank as $\hMat_{N}^{(\alpha)}(\lambda)$, completing the proof.
		\end{proof}

\end{document}
